\chapter{Data Science}

\section{Data Preparation}
\dfn{Missing Values}{
Handle via \texttt{fillna} (constant/forward/backward), mean/median, \texttt{interpolate}, \texttt{replace}, or \texttt{dropna}.
}
\dfn{Normalization vs Standardization}{
Rescale to [0,1] vs zero-mean/unit-variance; choose based on model sensitivity.
}
\dfn{Dummification}{
Create one-hot/indicator variables for categorical features.
}
\subsection*{Pandas Snippets}
\begin{lstlisting}
import pandas as pd, numpy as np
df.fillna(0, inplace=True)                      # constant
df.fillna(method="ffill", inplace=True)         # forward fill
df["Gender"].fillna("Unknown", inplace=True)    # categorical
df["x"].fillna(df["x"].mean(), inplace=True)    # mean impute
df.replace(np.nan, -99, inplace=True)           # sentinel
df.interpolate(method="linear", limit_direction="forward", inplace=True)
df.dropna(axis=0, how="any", inplace=True)      # drop rows with NA
pd.get_dummies(df, columns=["cat"])             # dummify
\end{lstlisting}

\section{Exploratory Data Analysis (EDA)}
\dfn{Content Identification}{
Understand columns, units, domains; inspect \texttt{dtypes}, uniques.
}
\dfn{Variable Types}{
Categorical vs numerical; continuous vs discrete.
}
\dfn{Granularity}{
Level of detail (row = transaction? day? customer?).
}
\dfn{Distribution Analysis}{
Summary stats + hist/kde/boxplots.
}
\dfn{Missingness Assessment}{
Counts/percents; visualize bars/heatmaps.
}
\dfn{Multivariate Analysis}{
Correlation, sparsity, cross-tabs, pairplots.
}
\subsection*{Quick Checks}
\begin{lstlisting}
data.dtypes
cat = data.select_dtypes(include="object")
for c in cat: print(c, data[c].unique())
data.isna().sum().plot.bar()
data.describe(include="all")
data.corr(numeric_only=True)
\end{lstlisting}

\section{Evaluation Metrics}
\dfn{Confusion Matrix}{
TP, TN, FP, FN for binary classification.
}
\dfn{Accuracy}{
\(\frac{TP+TN}{TP+FP+TN+FN}\); overall correct rate.
}
\dfn{Precision}{
\(\frac{TP}{TP+FP}\); focus when FP cost is high.
}
\dfn{Recall}{
\(\frac{TP}{TP+FN}\); focus when FN cost is high.
}
\dfn{Type I / II Errors}{
FP = Type I; FN = Type II.
}
\subsection*{Sklearn Example}
\begin{lstlisting}
from sklearn.metrics import confusion_matrix, precision_score, recall_score, accuracy_score
y_pred = model.predict(X_test)
print(confusion_matrix(y_test, y_pred))
print("precision", precision_score(y_test, y_pred))
print("recall", recall_score(y_test, y_pred))
print("accuracy", accuracy_score(y_test, y_pred))
\end{lstlisting}

\section{Process and Methodology}
\dfn{CRISP-DM / POST-DS}{
Phases: Business Understanding, Data Understanding, Data Preparation, Modelling, Evaluation, Deployment; POST-DS adds Organization, Scheduling, Tools.
}
\dfn{Process vs Organization vs Scheduling vs Tools}{
Separate concerns: workflow steps, roles/teams, timing, and technical stack.
}

\section{Reinforcement Learning (Intro)}
\dfn{Reinforcement Learning}{
Agent interacts with environment, takes actions, receives rewards, learns policy.
}
\dfn{Key Components}{
Agent, Environment, State, Action, Reward, Policy.
}
\dfn{Q-Learning}{
Learns \(Q(s,a)\) values to choose actions that maximize future reward.
}

\subsection*{Toy RL Loop (Q-learning skeleton)}
\begin{lstlisting}
import numpy as np
Q = np.zeros((n_states, n_actions))
alpha, gamma, eps = 0.1, 0.99, 0.1
for episode in range(500):
    s = env.reset()
    done = False
    while not done:
        a = np.argmax(Q[s]) if np.random.rand()>eps else np.random.randint(n_actions)
        s2, r, done, _ = env.step(a)
        Q[s,a] += alpha * (r + gamma*np.max(Q[s2]) - Q[s,a])
        s = s2
\end{lstlisting}

\section{Practice / Exam Prompts}
\pr{Imputation strategy}{
Given mixed numeric/categorical data with 5\% missing, propose imputation per type and justify when to drop rows/cols.
}
\pr{Granularity check}{
Explain how to detect granularity issues and why duplicates may indicate wrong grain.
}
\pr{Metric choice}{
Skin-cancer detection vs. spam filter: which metrics to prioritize (recall vs precision) and why.
}
\pr{Process mapping}{
Map a project to POST-DS: list tasks under Process, Organization, Scheduling, Tools.
}
\pr{RL components}{
Identify state/action/reward/policy in a trading-bot example.
}
